\chapter{TWO WHEELED INVERTED PENDULUM SIMULATION}
\label{ch:simulation}

\section{DISCRETIZATION}
\label{sec:discretization}

The linear model \eqref{eq:linss} is discretized at sample period $\Delta t$
using the matrix exponential,
\begin{equation}\label{eq:zoh}
F=e^{\mathcal{A}\Delta t},
\qquad
G=\left(\int_0^{\Delta t}e^{\mathcal{A}\tau}\,\dd\tau\right)\mathcal{B},
\end{equation}
which is exact for a zero-order hold input. The plant used for truth propagation
is the nonlinear model \eqref{eq:nl1}--\eqref{eq:nl2}, integrated at a step at
least an order of magnitude smaller than $\Delta t$, so that the estimator and
controller see a discretized measurement of a continuous truth rather than a
discretized truth.

\rev{\begin{remark}[The sample period enters the stability conditions]
The quantity $\sigmax(A)=\sigmax(F-\Km)$ depends on $\Delta t$ through
\eqref{eq:zoh}, so the admissible range of $\Gamma$ in
Corollary~\ref{cor:theta1} is a function of the sample rate. Reducing $\Delta t$
moves $F$ toward the identity and requires a correspondingly larger $\Km$ for
the same $\sigmax(A)$. The design must be rechecked whenever the loop rate
changes rather than assumed to carry over.
\end{remark}}

\section{SIMULATION MODEL}
\label{sec:simmodel}

\rev{The truth model used for the results of Chapter~\ref{ch:results} is the
nonlinear model of Section~\ref{sec:twipmodel} with three corrections. Each was
identified while reimplementing the simulations, and each can be checked
against material already in this thesis.

\paragraph{Back-EMF kinematics.} The motor stator is fixed to the pendulum
body, so the back-EMF and the friction of each motor depend on the rotation of
the wheel \emph{relative to the body}, not on the absolute wheel rate
$\dot\theta_w=\dot x/r$ of \eqref{eq:noslip}. In the upright-referenced
coordinates of \eqref{eq:linss} the wheel turns clockwise at $\dot x/r$ while
the body turns counterclockwise at $\dot\phi$, so the relative rate is
$\dot x/r+\dot\phi$, and the total motor torque is
\begin{equation}\label{eq:relemf}
T=\frac{2k_m}{R}\Bigl(V-k_e\bigl(\tfrac{\dot x}{r}+\dot\phi\bigr)\Bigr).
\end{equation}
The input terms of \eqref{eq:linss} already imply this kinematics: the torque
enters the wheel equation as $T/r$ and the pendulum equation as $+T$
(\eqref{eq:pendm}), which is the virtual work of $T$ acting through the relative
angle $x/r+\phi$. Using \eqref{eq:relemf} adds the tilt-rate terms
$\mathcal{A}_{24}=r\mathcal{A}_{22}$ and $\mathcal{A}_{44}=r\mathcal{A}_{42}$ to
\eqref{eq:linss}, which are zero in the original. It is also the form that makes the motor strictly dissipative at zero
voltage.

\paragraph{Back-EMF constant.} From the free-run data of Table~\ref{tab:motor},
the steady armature equation $V=Ri+k_e\omega$ at the free-run speed of
\SI{350}{rpm} (\SI{36.65}{rad/s}) gives
\[
k_e=\frac{\SI{12}{V}-(\SI{0.3}{A})(\SI{2.5}{\ohm})}{\SI{36.65}{rad/s}}
=\SI{0.307}{V.s/rad},
\]
about 85 times the value used in revision~7. The torque constant is retained.
The stall data give $k_m\approx\SI{0.155}{N.m/A}$, and the discrepancy
between the two is one reason Stage~0 of Appendix~\ref{app:validation} calls
for direct measurement.

\paragraph{Pendulum inertia.} The value of $I_p$ in Table~\ref{tab:phys} is
computed from a lumped-mass model of the plates that includes parallel-axis
terms, so it is the inertia about the wheel axis. Equations
\eqref{eq:nl1}--\eqref{eq:nl2} require the inertia about the center of
gravity, $I_p-M_pl^2=\SI{0.0169}{kg.m^2}$.

Table~\ref{tab:linvals} compares the resulting linearizations. The corrected
model has two to three orders of magnitude more back-EMF damping, a tilt input gain
$\mathcal{B}_4$ ten times larger, and a slightly slower unstable pole. The
right-half-plane zero of $x_1(s)/u(s)$, which matters in
Section~\ref{sec:cfinstab}, is present in both. Setting $x\equiv0$ in the
linearized equations shows that it lies at
$\sqrt{M_pgl/(I_p+M_pl^2\pm M_plr)}$, independent of the motor constants, with
the plus sign for the corrected model and the minus sign for revision~7. In the
corrected model it nearly coincides with the unstable pole, for a physical
reason. With the base still and zero voltage, the back-EMF torque is
$-(2k_mk_e/R)\dot\phi$, while holding the base still requires
$-M_plr\ddot\phi$. For a mode growing at rate $p$ the two agree when
$2k_mk_e/R=M_plrp$, and the corrected constants give \num{0.0283} and
\SI{0.0288}{N.m.s} at $p=\SI{5.95}{rad/s}$. The body therefore falls while the
base moves by less than 1\% of $r\phi$: the unstable mode barely appears in
$x_1$ and must be stabilized through the tilt.

\begin{table}[htbp]
\centering
\caption{Linearized model entries, revision 7 and corrected}
\label{tab:linvals}
\begin{tabular}{lcc}
\toprule
Quantity & Revision 7 & Corrected \\
\midrule
$\mathcal{A}_{22},\ \mathcal{A}_{24}$ & $-0.098,\ 0$ & $-12.65,\ -0.569$ \\
$\mathcal{A}_{42},\ \mathcal{A}_{43},\ \mathcal{A}_{44}$ & $-0.095,\ 39.7,\ 0$ & $-79.6,\ 57.0,\ -3.58$ \\
$\mathcal{B}_2,\ \mathcal{B}_4$ & $1.22,\ 1.18$ & $1.86,\ 11.67$ \\
Unstable pole (rad/s) & 6.30 & 5.953 \\
RHP zero of $x_1/u$ (rad/s) & 6.11 & 5.946 \\
\bottomrule
\end{tabular}
\end{table}

\paragraph{Sensors.} The accelerometer is modeled as the specific force at the
IMU location, \SI{2}{cm} above the wheel axis on the body, expressed in body
axes. The tilt presented to the estimators is therefore computed from a signal
that includes the robot's own acceleration, as on the hardware, rather than
being the true tilt plus white noise. The encoders measure the rotation of the
wheel relative to the body, so the encoder position is $x+r\phi$; it is
corrected by subtracting $r\hat\phi$ before use. Velocity is the difference of
successive quantized positions divided by $\Delta t$.}

\section{CONTROL}

The nominal control is the LQR law of Section~\ref{sec:lqr}, computed from the
nominal parameters of Table~\ref{tab:phys} and Table~\ref{tab:motor}. The extra
control of Section~\ref{sec:cffb} is added when the neural compensation is
active. The command filters \eqref{eq:cmdfilt} are integrated with the same
solver and step as the controller.

\section{UNCERTAINTY}

Uncertainty is introduced by propagating the truth model with perturbed physical
parameters while the observer and controller retain the nominal values. Because
the physical parameters enter only through the coefficients of the two
acceleration equations, this produces uncertainty in those two equations and
nowhere else, which is why the selection matrix is
\begin{equation}\label{eq:Bmso}
\Bm=\begin{bmatrix}0&1&0&0\\0&0&0&1\end{bmatrix}^T .
\end{equation}
The uncertainty vector has dimension four in the full state, but only the second
and fourth entries are assigned, so $p=2$ and $\Bm\in\R^{4\times2}$ satisfies
Assumption~\ref{as:B}. Unmodeled dynamics are introduced separately by retaining
in the truth model the nonlinear terms that the linearization
\eqref{eq:smallangle}--\eqref{eq:smallrate} discards.

\section{SIMULATED NOISE}
\label{sec:noise}

Sensor noise is not white in general, but both the gyroscope and the
accelerometer can be modeled with three principal sources of error
\cite{gebreegziabher2004design,flenniken2005modeling}. Both sensors are modeled as
\begin{equation}\label{eq:sensormodel}
m=m_r+c_r+b_r+w_r ,
\end{equation}
where $m$ is the sensor reading, $m_r$ the true measurement, $c_r$ a constant
bias that changes on each initialization of the sensor, $b_r$ the walking bias,
and $w_r$ a white noise process with variance $\sigma_r^2$.

The constant bias is easily remedied and is removed on initialization of the
sensor in digital logic, as described in Section~\ref{sec:measurements}. The
other two remain and must be considered in simulation. White noise is the
dominant source in low cost sensors and is characterized by finding the variance
of a long run of data taken with the sensor stationary. The walking bias is more
involved but can be modeled as a first-order Markov process,
\begin{equation}\label{eq:walkingbias}
\dot b_r=-\frac{1}{\tau_r}b_r+w_{b_r},
\end{equation}
where $\tau_r$ is the time constant of the Markov process and $w_{b_r}$ is white
noise with intensity $2\sigma_{b_r}^2/\tau_r$, so that $\sigma_{b_r}^2$ is the
steady-state variance of $b_r$. The walking bias captures the tendency of
the sensor bias to drift over time. Its parameters are obtained by performing an
autocorrelation analysis of a long run of filtered data; further detail is given
by Flenniken \cite{flenniken2005modeling}.

Allan variance plots are the standard method of characterizing the accuracy of
gyroscopes and accelerometers. Figure~\ref{fig:allan} shows an Allan variance
plot of measured gyroscope data against data simulated from
\eqref{eq:sensormodel}--\eqref{eq:walkingbias}. The agreement is not perfect and
other error sources are clearly present, but the result is far more
representative than simply adding white noise to the measurements.

The resulting parameters, measured from the sensors used in the TWIP
implementation, are given in Table~\ref{tab:noise}. Encoder measurements carry
quantization error set by the counts per revolution and the gear ratio; the
velocity quantization is the position quantization divided by $\Delta t$,
because velocity is obtained by differencing successive position readings.

\begin{table}[htbp]
\centering
\caption{Noise simulation parameters}
\label{tab:noise}
\begin{tabular}{lcc}
\toprule
Parameter & Tilt angle & Gyroscope \\
\midrule
$\sigma_r^2$    & $4.0474\times10^{-6}$ \si{rad^2}
                & $0.0012$ (\si{\degree/s})$^2$ \\
$\sigma_{br}^2$ & $1.8929\times10^{-7}$ \si{rad^2}
                & $1.0986\times10^{-4}$ (\si{\degree/s})$^2$ \\
$\tau_r$        & \SI{2160}{s} & \SI{1610}{s} \\
\midrule
\multicolumn{2}{l}{Encoder position quantization}
                & \SI{1.47e-4}{m} \\
\multicolumn{2}{l}{Encoder velocity quantization}
                & $1.47\times10^{-4}/\Delta t$ \si{m/s} \\
\bottomrule
\end{tabular}
\end{table}

\begin{figure}[htbp]
\centering
\includegraphics[width=0.75\textwidth]{allan_variance}
\caption{Gyroscope Allan deviation: measured (1.95 million samples at
\SI{140}{Hz}, stationary) against data simulated from
\eqref{eq:sensormodel}--\eqref{eq:walkingbias} with the parameters of
Table~\ref{tab:noise}, and against white noise alone.}
\label{fig:allan}
\end{figure}

\rev{\begin{remark}[The sample rate trade]\label{rem:dttrade}
The encoder velocity quantization interacts with the sample rate in the opposite
direction from everything else. Halving $\Delta t$ improves the fidelity of
\eqref{eq:zoh} but doubles the velocity quantization noise. There is therefore
an interior optimum in $\Delta t$ for a fixed encoder resolution, and raising the
loop rate without raising the encoder resolution eventually degrades the
velocity estimate. A procedure for locating this optimum from logged data is
given in Appendix~\ref{app:validation}.
\end{remark}}

\section{NONLINEARITIES}

The Pololu DC motors exhibit several obvious nonlinearities that can be
accounted for in simulation. Friction is an important nonlinearity to consider,
but owing to the difficulty of evaluating the forces and obtaining an accurate
model it was left out of the simulation; its effect is absorbed into $d_1$ and
$d_2$ and forms part of what the networks estimate.

\subsection{Saturation}

The commanded voltage is limited by the supply rail,
\begin{equation}\label{eq:sat}
\mathrm{sat}(u)=
\begin{cases}
u_{\max}, & u>u_{\max},\\
u, & \abs{u}\le u_{\max},\\
-u_{\max}, & u<-u_{\max}.
\end{cases}
\end{equation}
Saturation is applied to the control signal output in all simulation cases.

\subsection{Deadzone}

Below a threshold voltage the motor does not turn. The standard deadzone model
is continuous at the breakaway points, but the physical motors exhibit a jump
discontinuity: the output steps to a finite value at breakaway rather than
rising continuously from zero. The model used is
\begin{equation}\label{eq:deadzone}
\mathcal{D}(u)=
\begin{cases}
m_d\bigl(u-\delta_+\bigr)+j_+, & u>\delta_+,\\
0, & -\delta_-\le u\le\delta_+,\\
m_d\bigl(u+\delta_-\bigr)-j_-, & u<-\delta_- ,
\end{cases}
\end{equation}
with $\delta_\pm$ the breakaway voltages, $m_d$ the slope outside the deadband,
and $j_\pm$ the jump magnitudes.

\begin{figure}[htbp]
\centering
\includegraphics[width=0.55\textwidth]{deadzone}
\caption{Deadzone nonlinearity with jump discontinuities.}
\label{fig:deadzone}
\end{figure}

\subsection{Backlash}

Gear backlash is modeled as hysteresis \cite{tao1993backlash,durdevic2013hybrid} in which the output holds while the input
reverses through a deadband of half-width $\Delta_b$,
\begin{equation}\label{eq:backlash}
\dot y=
\begin{cases}
m_b\dot u, & \dot u>0 \text{ and } y=m_b\bigl(u-\Delta_b\bigr),\\
m_b\dot u, & \dot u<0 \text{ and } y=m_b\bigl(u+\Delta_b\bigr),\\
0, & \text{otherwise.}
\end{cases}
\end{equation}

\begin{figure}[htbp]
\centering
\includegraphics[width=0.55\textwidth]{backlash}
\caption{Backlash nonlinearity: response to a decaying sinusoidal input with
$m_b=1$, $\Delta_b=0.5$.}
\label{fig:backlash}
\end{figure}

\rev{\begin{remark}[Backlash is outside the approximation hypothesis]
\label{rem:backlash}
On the physical robot the backlash was measured at less than one degree of
output-shaft rotation, although the recorded runs show a stick-slip limit cycle
near upright (Section~\ref{sec:momtest}). Two observations follow. First, because the encoders are
mounted on the motor shaft, this deadband is invisible to the measurement, so it
appears to the estimator as an unmodeled input nonlinearity rather than as a
measurable state discrepancy. Second, \eqref{eq:backlash} is not smooth and
therefore violates the smoothness assumption underlying
Assumption~\ref{as:nn}. The networks can approximate it only in an averaged
sense, and the residual appears in $\epsilon_N$ and in $w_{i0}$. This is the
correct explanation for the persistent oscillation in the uncertainty estimates
reported in Chapter~\ref{ch:results}: the limit cycle is the signature of a
nonsmooth input nonlinearity that no smooth approximator can cancel, not a
tuning failure.
\end{remark}}

\section{DISCRETE KALMAN FILTER}
\label{sec:dkf}

The discrete Kalman filter is used for comparison against the DMSO. The Kalman
filter, developed by Kalman in 1960, provides the statistically optimal state
estimate for a linear system with measurements corrupted by Gaussian,
zero-mean, uncorrelated, white noise \cite{kalman1960}. The filter processes
measurements together with knowledge of the system dynamics to produce a more
accurate state estimate. The version used here is the discrete-time filter and
follows Simon \cite{simon2006optimal}.

Assume the dynamic system is given by
\begin{equation}\label{eq:kfplant}
x_{k+1}=Fx_k+Gu_k+w_k ,
\end{equation}
with measurements
\begin{equation}\label{eq:kfmeas}
y_k=Hx_k+v_k ,
\end{equation}
where $w_k$ and $v_k$ are zero-mean, uncorrelated, white Gaussian sequences with
\begin{equation}\label{eq:kfstats}
w_k\sim\mathcal{N}\bigl(0,Q_{\mathrm{K}}\bigr),
\qquad
v_k\sim\mathcal{N}\bigl(0,R_{\mathrm{K}}\bigr).
\end{equation}
$Q_{\mathrm{K}}$ is the process noise covariance and $R_{\mathrm{K}}$ the
measurement noise covariance; both may be adjusted for performance. States with
greater uncertainty in the dynamics can be assigned larger values in
$Q_{\mathrm{K}}$, which causes the filter to weight the measurements more
heavily. $R_{\mathrm{K}}$ can be obtained accurately by statistical analysis of
the available measurements.

The filter is given by the update equations
\begin{align}
P_{k+1}^-&=FP_k^+F^T+Q_{\mathrm{K}},\label{eq:kf1}\\[2pt]
L_{k+1}&=P_{k+1}^-H^T\bigl(HP_{k+1}^-H^T+R_{\mathrm{K}}\bigr)^{-1},\label{eq:kf2}\\[2pt]
\xhat_{k+1}^-&=F\xhat_k^++Gu_k ,\label{eq:kf3}\\[2pt]
\xhat_{k+1}^+&=\xhat_{k+1}^-+L_{k+1}\bigl(y_{k+1}-H\xhat_{k+1}^-\bigr),\label{eq:kf4}\\[2pt]
P_{k+1}^+&=\bigl(I-L_{k+1}H\bigr)P_{k+1}^-\bigl(I-L_{k+1}H\bigr)^T
+L_{k+1}R_{\mathrm{K}}L_{k+1}^T ,\label{eq:kf5}
\end{align}
where the superscripts $-$ and $+$ denote a priori and a posteriori quantities,
$P$ is the estimation error covariance, $L$ is the Kalman gain, and $\xhat$ the
state estimate. The equations are processed sequentially from \eqref{eq:kf1} to
\eqref{eq:kf5}. Equation \eqref{eq:kf5} is the Joseph form, which preserves
symmetry and positive definiteness of the covariance under finite-precision
arithmetic. Measurements need not arrive at the same rate as the dynamics are
propagated; \eqref{eq:kf4} may be processed only at intervals where measurements
are available. In the simulations here the measurements are assumed available at
the dynamics rate.

\rev{\subsection{Three Baselines}

A comparison between the DMSO and a Kalman filter must be constructed carefully
to be meaningful. The Kalman filter is optimal for a linear system driven by
zero-mean white process noise of known covariance. It is not designed to reject
a deterministic modeling error, and demonstrating that it fails to do so
establishes little about the DMSO. Three baselines are therefore used in
Chapter~\ref{ch:results}:
\begin{enumerate}[label=(\alph*),leftmargin=*]
\item a Kalman filter with $Q_{\mathrm{K}}$ tuned to the sensor noise alone,
which is the naive baseline and the one used in revision~7;
\item a Kalman filter with $Q_{\mathrm{K}}$ inflated to account for the
parameter error, which is what a competent practitioner would actually do;
\item an augmented-state Kalman filter in which the uncertainty is modeled as a
random-walk state and estimated alongside the dynamic states, which is the
closest linear analogue of the DMSO and the fairest comparison.
\end{enumerate}
For baseline (c) the state is augmented as
$x^{\mathrm{aug}}=\begin{bmatrix}x^T&f^T\end{bmatrix}^T$ with
\begin{equation}\label{eq:augkf}
F^{\mathrm{aug}}=\begin{bmatrix}F&\Bm\\0&I_p\end{bmatrix},
\qquad
G^{\mathrm{aug}}=\begin{bmatrix}G\\0\end{bmatrix},
\qquad
H^{\mathrm{aug}}=\begin{bmatrix}H&0\end{bmatrix},
\end{equation}
and the process noise covariance on the appended block set to the assumed
random-walk intensity of the uncertainty. Claims of improvement are stated
against (c).}
